% Author: Grayson Orr
% Course: ID721001: Mobile Application Development

\documentclass{article}
\author{}

\usepackage{graphicx}
\usepackage{wrapfig}
\usepackage{enumerate}
\usepackage{hyperref}
\usepackage[margin = 2.25cm]{geometry}
\usepackage[table]{xcolor}
\usepackage{fancyhdr}
\hypersetup{
  colorlinks = true,
  urlcolor = blue
}
\setlength\parindent{0pt}
\pagestyle{fancy}
\fancyhf{}
\rhead{College of Engineering, Construction \& Living Sciences\\Bachelor of Information Technology}
\lfoot{Project\\Version 2, Semester Two, 2023}
\rfoot{\thepage}
 
\begin{document}

\begin{figure}
	\centering
	\includegraphics[width=50mm]{../img/logo.png}
\end{figure}

\title{College of Engineering, Construction \& Living Sciences\\Bachelor of Information Technology\\ID721001: Mobile Application Development\\Level 7, Credits 15\\\textbf{Project}}
\date{}
\maketitle

\section*{Assessment Overview}
In this \textbf{individual} assessment, you will develop a mobile application using \textbf{React Native} \& \textbf{Expo}, \& publish it to \textbf{Google Play Store} or \textbf{Apple App Store}. Also, you will provide documentation that addresses several aspects of the mobile application development process. In addition to the mobile application \& documentation, you will present the mobile application and answer follow up questions via a video recording.

\section*{Learning Outcomes}
At the successful completion of this course, learners will be able to:
\begin{enumerate}
	\item Implement \& publish complete, non-trivial, industry-standard mobile applications following sound architectural \& code-quality standards.
	\item Identify relevant use cases for a mobile computing scenario \& incorporate them into an effective user experience design.
	\item Follow industry standard software engineering practice in the design of mobile applications.
\end{enumerate}

\section*{Assessments}
\renewcommand{\arraystretch}{1.5}
\begin{tabular}{|c|c|c|c|}
	\hline
	\textbf{Assessment} & \textbf{Weight} & \textbf{Due Date}    & \textbf{Learning Outcomes} \\ \hline
	Project: Client Mobile Application          & 40\%            & 10-11-2023 (Friday at 4.59 PM)  & 1, 2, 3                    \\ \hline
	Practical: Skills-Based           & 20\%            & 22-09-2023 (Friday at 4.59 PM)  & 1, 2, 3                    \\ \hline
	Presentation       & 20\%            & 10-11-2023 (Friday at 4.59 PM) & 2, 3                       \\ \hline
\end{tabular}

\section*{Conditions of Assessment}
You will complete majority of this assessment during your learner-managed time. However, there will be time during class to discuss the requirements and your progress on this assessment. This assessment will need to be completed by \textbf{Friday, 10 November 2023} at \textbf{4.59 PM}.

\section*{Pass Criteria}
This assessment is criterion-referenced (CRA) with a cumulative pass mark of \textbf{50\%} over all assessments in \textbf{ID721001: Mobile Application Development}.

\section*{Authenticity}
All parts of your submitted assessment \textbf{must} be completely your work. Do your best to complete this assessment without using an \textbf{AI generative tool}. You need to demonstrate to the course lecturer that you can meet the learning outcome for this assessment. \\
 
 However, if you get stuck, you can use an \textbf{AI generative tool} to help you get unstuck, permitting you to acknowledge that you have used it. In the assessment's repository \textbf{README.md} file, please include what prompt(s) you provided to the \textbf{AI generative tool} and how you used the response(s) to help you with your work. It also applies to code snippets retrieved from \textbf{StackOverflow} and \textbf{GitHub}. \\
 
 Failure to do this may result in a mark of \textbf{zero} for this assessment.

\section*{Policy on Submissions, Extensions, Resubmissions \& Resits}
The school's process concerning submissions, extensions, resubmissions \& resits complies with \textbf{Otago Polytechnic} policies. Learners can view policies on the \textbf{Otago Polytechnic} website located at \href{https://www.op.ac.nz/about-us/governance-and-management/policies}{https://www.op.ac.nz/about-us/governance-and-management/policies}.

\section*{Submission}
You \textbf{must} submit all project files via \textbf{GitHub Classroom}. Here is the URL to the repository you will use for your submission – \href{https://classroom.github.com/a/yitUo0I6}{https://classroom.github.com/a/yitUo0I6}.  Create a \textbf{.gitignore} \& add the ignored files in this resource - \href{https://raw.githubusercontent.com/github/gitignore/main/Node.gitignore}{https://raw.githubusercontent.com/github/gitignore/main/Node.gitignore}. The latest project files in the \textbf{master} or \textbf{main} branch will be used to mark against the \textbf{Functionality} criterion. Please test before you submit. Partial marks \textbf{will not} be given for incomplete functionality. Late submissions will incur a \textbf{10\% penalty per day}, rolling over at \textbf{5:00 PM}.

\section*{Extensions}
Familiarise yourself with the assessment due date. If you need an extension, contact the course lecturer before the due date. If you require more than a week's extension, a medical certificate or support letter from your manager may be needed.

\section*{Resubmissions}
Learners may be requested to resubmit an assessment following a rework of part/s of the original assessment. Resubmissions are to be completed within a negotiable short time frame \& usually \textbf{must} be completed within the timing of the course to which the assessment relates. Resubmissions will be available to learners who have made a genuine attempt at the first assessment opportunity \& achieved a \textbf{D grade (40-49\%)}. The maximum grade awarded for resubmission will be \textbf{C-}.

\section*{Resits}
Resits \& reassessments \textbf{are not} applicable in \textbf{ID721001: Mobile Application Development}.

\section*{Instructions}
You will need to submit a mobile application \& documentation that meet the following requirements:

\subsection*{Functionality - Learning Outcomes 1, 2, 3 (60\%)}
\begin{itemize}
	\item The mobile application needs to run without code or file structure modification in \textbf{Visual Studio Code}.
	\item Adhere to the functionality requirements outlined by the client.
	\item Usable on a variety of mobile devices, i.e., devices with different screen sizes.
	\item Free of bugs that significantly effect the usability.
	\item The mobile application is published to \textbf{Google Play Store} or \textbf{Apple App Store}.
	      \begin{itemize}
		      \item To published to \textbf{Google Play Store} or \textbf{Apple App Store}, you will need an account. The account's credentials will be privately given to you on \textbf{Microsoft Teams}. \textbf{Do not} disable any applications published on this account.
	      \end{itemize}
	\item Ability to download the mobile application from \textbf{Google Play Store} or \textbf{Apple App Store} on to a variety of mobile devices.
\end{itemize}

\subsection*{Code Elegance - Learning Outcomes 1, 3 (45\%)}
\begin{itemize}
	\item \textbf{Kotlin} \& \textbf{XML} files contain no magic numbers/strings. Store the values in the appropriate \textbf{XML} files. For example, numbers \textbf{must} be stored in an \textbf{integer.xml} or a \textbf{dimens.xml} file \& strings \textbf{must} be stored in a \textbf{strings.xml}.
	\item Idiomatic use of control flow, data structures \& in-built functions.
	\item Code adheres to \textbf{DRY}, \textbf{KISS} \& the \textbf{Model-View-ViewModel} architectural pattern.
	\item Adhere to an \textbf{OO} architecture. For example, classes, functions, concise naming \& functions assigned to the correct classes.
	\item Efficient algorithmic approach.
	\item Commented code is documented using \textbf{KDoc}. The purpose of each class and function \textbf{must} be explained.
	\item API keys are stored \& retrieved from \textbf{local.properties}.
	\item \textbf{Kotlin} \& \textbf{XML} files are code formatted.
	\item No unused code \& resources.
\end{itemize}

\subsection*{Documentation - Learning Outcomes 2, 3 (20\%)}
\begin{itemize}
	\item Provide the following in your repository \textbf{README.md} file:
	\begin{itemize}
		\item Game overview:
		\begin{itemize}
			\item What is the game's concept and genre?
		\end{itemize}
		\item Gameplay mechanics:
		\begin{itemize}
			\item What is the game's core gameplay loop?
			\item How do the players interact with the game and control the characters or objects?
			\item 
		\end{itemize}
		\item Art and visuals:
		\begin{itemize}
			\item What art style is used?
		\end{itemize}
		\item Audio and sound:
		\begin{itemize}
			\item What audio and sound effects are used?
		\end{itemize}
		\item If applicable, known bugs.
	\end{itemize}
\end{itemize}

\subsection*{Presentation - Learning Outcomes 2, 3 (20\%)}
\begin{itemize}
	\item Present the mobile application via a video recording. In addition, you need to answer the following:
	\begin{itemize}
		\item How did you plan and prioritise features throughout the development process?
		\item What tools and technologies did you utilise to streamline your development workflow?
		\item How did you handle potential challenges, such as time management and motivation?
		\item What strategies did you employ to maintain code quality and avoid technical debt during the development process?
		\item How did you handle testing and debugging of the mobile application?
	\end{itemize}
\end{itemize}

\subsection*{Additional Information}
\begin{itemize}
	\item \textbf{Do not} rewrite your \textbf{Git} history. It is important that the course lecturer can see how you worked on your assessment over time.
\end{itemize}

\end{document}